\section{Présentation de la turbine Pelton}
%Elle doit contenir les plages d’utilisation de la turbine,
%la classification de la turbine, l’explication sommaire de son fonctionnement et le parcours de l’eau
%dans les différentes composantes (nommées correctement).
La turbine du laboratoire est une turbomachine à impulsion et à axe horizontal munie d'une roue "Pelton 18-G", soit une roue de 381 mm de diamètre, de géométrie fixe et constituée d'agets.
Son fabricant est l'entreprise américaine "The Pelton water wheel company".
Cette roue est mise en mouvement par un jet d'eau tangeant  qui vient heurter les augets de la roue et qui provient d'un injecteur dont l'ouverture peut être contrôllée par un vérin affublé d'un pointeau. L'eau est acheminée vers l'injecteur par une canalisation d'amenée prenant sa source au sous-sol du laboratoire, où une pompe impose une hauteur de chute artificielle à l'eau d'un réservoir d'alimentation.
Après avoir percuté les augets, l'eau est devié de 160 à 170 degrés avant d'être retournée au réservoir par le biais d'une seconde conduite sous la roue. 
La roue de la turbine Pelton tourne donc grâce au transfert de quantité de mouvement entre le jet et les augets, le premier obtenant son énergie cinétique par la conversion de l'énergie potentielle de la hauteur de chute.
Contrairement au transfert d'énergie dans une turbine à réaction, celui dans une turbine Pelton s'opère à pression atmosphérique.  

\begin{table}[H]
\centering
\caption{Plages d'opérations de la turbine}
    \label{tab:Op_table}
    \begin{tabular}{|c|c|}
        \hline
	\textbf{Puissance nominale} & 7 kW \\
        \hline
	\textbf{Débit nominal} & 0.020 $m^3/s$ \\
        \hline
	\textbf{Chute nominale} & 38.76 m \\
        \hline
	\textbf{Vitesse nominale} & 625 tr/min \\
        \hline
	\textbf{Plage de vitesse de rotation} & 550 à 830 tr/min \\
        \hline
	\textbf{Diamètre de la roue} & 0.381 m \\
        \hline
	\textbf{Ouverture du vérin de l'injecteur} & 0 à 40 mm \\
        \hline
	\textbf{Angle de augets} & 160 degrés \\
        \hline
   \end{tabular}
\end{table}	
